\section{Introducción: los desafíos de ejecutar Agentes durante periodos prolongados}

A medida que las capacidades de los Agentes de IA se fortalecen, los desarrolladores comienzan a exigirles que asuman tareas complejas que abarcan horas e incluso días. Sin embargo, mantener un progreso coherente entre múltiples ventanas de contexto sigue siendo un desafío de ingeniería abierto.

\begin{knowledgebox}{¿Qué es una ventana de contexto?}
Una ventana de contexto es la longitud total de texto que un modelo de lenguaje grande puede "percibir" en una única llamada de inferencia. Incluso modelos avanzados como Claude tienen una ventana de contexto limitada. Cuando una tarea requiere múltiples llamadas (varios sesiones) para completarse, cada nueva sesión comienza desde cero —sin memoria de la ronda anterior—. Este es el problema central que enfrentan los Agentes de ejecución prolongada.
\end{knowledgebox}

El artículo de esta entrada del equipo de ingeniería de Anthropic propone una metáfora ilustrativa: imagine un proyecto de software completado por turnos por varios ingenieros, pero donde cada ingeniero no recuerda absolutamente nada de lo ocurrido en el turno anterior al llegar. Esta es exactamente la situación a la que se enfrentan los Agentes de ejecución prolongada: una ventana de contexto limitada frente a proyectos complejos que no pueden caber dentro de una sola ventana.

Para abordar este problema, Anthropic desarrolló un enfoque de dos etapas basado en el SDK de Agentes de Claude:
\begin{enumerate}
    \item \textbf{Initializador Agente}: configura el entorno durante la primera ejecución;
    \item \textbf{Agente de Codificación}: avanza progresivamente la tarea en cada sesión subsiguiente y deja un registro de trabajo claro para la próxima sesión.
\end{enumerate}

\begin{importantbox}{Idea central: un relevo como lo haría un ingeniero humano}
La clave de este enfoque es permitir que los Agentes trabajen como ingenieros de software humanos excelentes —redactando documentación de relevo, cometiendo código y registrando el progreso antes de cada cambio de turno—. Mediante la combinación del archivo \texttt{claude-progress.txt} con el historial de git, las nuevas sesiones de Agente pueden comprender rápidamente el estado actual del proyecto.
\end{importantbox}

\subsection{Resumen de la sección}

El desafío central de los Agentes de ejecución prolongada radica en la pérdida de estado entre sesiones debido a la limitación de la ventana de contexto. La solución de Anthropic adopta una arquitectura de dos etapas de "inicialización más codificación progresiva", que cierra las brechas entre sesiones mediante una gestión de entorno estructurada y un seguimiento del progreso.

